\documentclass[embeddedlogo]{../ufsc-thesis-rn46-2019}

\usepackage[T1]{fontenc} % fontes
\usepackage[utf8]{inputenc} % UTF-8
\usepackage{lipsum} % Gerador de texto
\usepackage{pdfpages} % Inclui PDF externo (ficha catalográfica)

% Usado para mostrar código
\usepackage[cache=false]{minted}
\newmintinline[mt]{latex}{fontsize=\normalsize}
\newmintinline[mft]{latex}{fontsize=\footnotesize}
\setminted{fontsize=\tiny,linenos,xleftmargin=2em}
\setmintedinline{breaklines,breakbytokenanywhere}

%%%%%%%%%%%%%%%%%%%%%%%%%%%%%%%%%%%%%%%%%%%%%%%%%%%%%%%%%%%%%%%%%%%%
%%% Configurações da classe (dados do trabalho)                  %%%
%%%%%%%%%%%%%%%%%%%%%%%%%%%%%%%%%%%%%%%%%%%%%%%%%%%%%%%%%%%%%%%%%%%%

% Informações para capa e folha de rosto/certificacao
\titulo{Template \LaTeX{} seguindo a RN 46/2019/CPG da UFSC}
\autor{Omar Ravenhurst}
\data{1 de agosto de 2019} % ou \today
\instituicao{Universidade Federal de Santa Catarina}
\centro{Centro Tecnológico}
\local{Florianópolis} % Apenas cidade! Sem estado
\programa{Programa de Pós-Graduação em Ciência da Computação}
% Os dois próximos itens são usados para gerar o \preambulo
\tese % ou \dissertacao
\titulode{Doutor em Ciência da Computação}
\orientador{Prof. Dr. Ben Trovato}
\coorientador{Prof. Dr. Lars Thørväld}

% Membros da banca e coordenador
% As regras da BU agora exigem que Dr. apareça depois do nome
\membrobanca{Prof. Valerie Béranger, Dr.}{Universidade Federal de Santa Catarina}
\membrobanca{Prof. Mordecai Malignatus, Dr.}{Universidade Federal de Santa Catarina}
\membrobanca{Prof. Huifen Chan, Dr.}{Universidade Federal de Santa Catarina}
% Atenção! o template da BU e o documento que apresenta as regras continua
% usando Dr antes do nome para Orientador e Coordenador!
\coordenador{Prof. Dr. Charles Palmer}

\begin{document}

%%%%%%%%%%%%%%%%%%%%%%%%%%%%%%%%%%%%%%%%%%%%%%%%%%%%%%%%%%%%%%%%%%%%
%%% Principais elementos pré-textuais                            %%%
%%%%%%%%%%%%%%%%%%%%%%%%%%%%%%%%%%%%%%%%%%%%%%%%%%%%%%%%%%%%%%%%%%%%

% Inicia parte pré-textual do documento capa, folha de rosto, folha de
% aprovação, aprovação, resumo, lista de tabelas, lista de figuras, etc.
\pretextual%
\imprimircapa%
\imprimirfolhaderosto*
% Atenção! \cleardoublepages são inseridos automaticamente
% Atenção! esse \protect é importante
\protect\incluirfichacatalografica{ficha.pdf}
\imprimirfolhadecertificacao
\listoffigures*  % O * evita que apareça no sumário
\tableofcontents*%

%%%%%%%%%%%%%%%%%%%%%%%%%%%%%%%%%%%%%%%%%%%%%%%%%%%%%%%%%%%%%%%%%%%%
%%% Corpo do texto                                               %%%
%%%%%%%%%%%%%%%%%%%%%%%%%%%%%%%%%%%%%%%%%%%%%%%%%%%%%%%%%%%%%%%%%%%%
\textual%

\chapter{Introdução}

Bem-vindo ao guia de usuário da classe \texttt{ufsc-thesis-rn46-2019}. Essa
classe é um conjunto de customizações aplicadas à classe
\href{https://ctan.org/pkg/abntex2}{\abnTeX} e ao pacote \texttt{abntex2cite}.
O objetivo da classe \texttt{ufsc-thesis-rn46-2019} é simplório: adequar o
\abnTeX{} às \href{http://portal.bu.ufsc.br/normalizacao/}{normas emitidas pela
Biblioteca Universitária da UFSC} em sequência à
\href{https://repositorio.ufsc.br/handle/123456789/197121}{Resolução Normativa
nº 46/2019/CPG}.


\section{Autores, suporte e atualizações}

Essa classe foi escrita inicialmente por dois alunos do
\href{http://ppgcc.posgrad.ufsc.br/}{PPGCC da UFSC}:
\href{mailto:alexishuf@gmail.com}{Alexis Huf} e
\href{mailto:gustavo.zambonin@posgrad.ufsc.br}{Gustavo Zambonin}.  Há o risco de
esse arquivo não ser atualizado a cada \textit{pull request}, então confira a
lista de mártires no GitHub. Essa classe é mantida no repositório
\href{https://github.com/alexishuf/ufsc-thesis-rn46-2019/}{alexishuf/ufsc-thesis-rn46-2019}.
Atualizações podem ser encontradas nesse repositório. \textit{Issues} e PRs são
bem vindos.

\subsection{Registro de mudanças}

Lista de versões (pelo menos das versões que receberam um número):
\begin{description}
\item[v0.1] (2019-08-02) Primeira versão completa e documentada.
\item[v0.1-alpha] (2019-08-01) Primeira versão para ajudar alunos com pouco
    prazo de entrega.
\end{description}

\section{Guia rápido}
\label{sec:quick}

A classe \texttt{ufsc-thesis-rn46-2019} deveria ser encarada como um
\emph{drop-in} para o \abnTeX. Você pode começar a escrever uma tese do zero
usando apenas esta classe, mas talvez você queira usar algum template indicado
por algum colega. Na maioria dos casos, bastará usar
\mt|\documentclass{ufsc-thesis-rn46-2019}|, ajustando o caminho se necessário.
Há duas principais formas de incluir essa classe no seu projeto.

\begin{enumerate}
  \item Copie o arquivo \texttt{ufsc-thesis-rn46-2019.cls}. Sim, isso funciona
      pois o logo da UFSC está embutido em base64 dentro da classe. Para que
        isso funcione, você deve compilar com \texttt{-shell-escape} em um
        ambiente \textit{UNIX-like}, como no arquivo \texttt{Makefile} de
        exemplo.
  \item Copie a pasta do repositório do GitHub para dentro do seu projeto.
      Nesse caso você deverá adicionar o caminho dessa pasta ao fazer o
        \mt|\documentclass|. Há 3 formas de incluir essa pasta:
  \begin{enumerate}
    \item Use a opção ``\emph{Download ZIP}'' do GitHub se você não sabe o que
        é Git.
    \item Via \texttt{git submodule}, se você já está usando Git.
    \item Via \texttt{git clone} se não estiver usando Git no seu projeto
        (\emph{shame on you}) ou se estiver usando no Overleaf, onde submódulos
          não funcionam.
    \end{enumerate}
\end{enumerate}

Esse guia rápido assume que você está começando do zero, sem um template
anterior, com intuito de ser didático. Se o seu trabalho for escrito em língua
inglesa, você deverá passar uma opção \mt|english| para a classe. Caso contrário
não precisa passar nenhuma opção. Logo após chamar a classe, você deverá
fornecer dados para a classe (e para o \abnTeX). Veja um preâmbulo completo de
exemplo na \autoref{fig:preambulo}. Os próximos parágrafos. Segue um rápido
\emph{overview} sobre o significado de cada um dos comandos usados (a maior
parte deles continua funcionando como no \abnTeX).

\begin{itemize}
  \item \mt|\titulode|, \mt|\autor|, \mt|\instituicao|, \mt|\orientador| e
      \mt|\coorientador|: esses comandos continuam com o mesmo significado e
        uso que possuem no \abnTeX. Para facilitar o uso, também foram
        adicionados \mt|\orientadora| e \mt|\coorientadora|.
  \item \mt|\data|: as regras da UFSC exigem que apenas o ano esteja presente
      na capa e folha de rosto. No entanto, a nova folha de certificação, que é
        gerada por essa classe precisa da data completa. Forneça a data
        completa \textbf{em português} (mesmo para documentos em inglês). A
        classe irá extrair o ano.
  \item \mt|\programa| e \mt|\centro|: nome do Programa de Pós-Graduação, por
      extenso, e nome do centro, \emph{e.g.} Centro Tecnológico.
  \item \mt|\tese| e \mt|\dissertacao|: o texto a ser colocado abaixo do título
      na folha de rosto tem regras bem definidas. Você deve usar um desses dois
        comandos para indicar o tipo de trabalho (e consequentemente o nível).
        Evidentemente, são mutuamente exclusivos.
  \item \mt|\preambulo|: Esse comando fornecido pelo \abnTeX{} é desnecessário
      nessa classe. Entretanto, você pode usar ele para sobrescrever o texto
        gerado automaticamente a partir do uso de \mt|\tese| e \mt|\titulode|.
\end{itemize}

A ``folha de certificação'', que substituiu a antiga folha de aprovação, agora
é gerada em \LaTeX. O nome do orientador já foi fornecido anteriormente, resta
apenas indicar os membros avaliadores da banca e o coordenador do programa:
\begin{itemize}
  \item \mt|\membrobanca{<nome>}{<instituição>}|: adiciona um membro da banca.
      \textbf{Atenção}: para membros avaliadores, o ``Dr.'' deve ser inserido
        após o nome e o ``Prof(a).'' deve preceder o nome.
  \item \mt|\coordenador| e \mt|\coordenadora|: Configura o nome do(a)
      coordenador(a) do Programa de Pós-Graduação.
\end{itemize}

\begin{figure}[tb]
  \centering
  \caption{Preâmbulo de uma tese (ou dissertação) típica usando esta classe.}
  \label{fig:preambulo}

  \begin{minted}{latex}
\documentclass[english]{ufsc-thesis-rn46-2019/ufsc-thesis-rn46-2019}

% \usepackage's a gosto

\titulo{Template \LaTeX~ seguindo a RN 46/2019/CPG da UFSC}
\autor{Fulano da Silva}
\data{1 de Agosto de 2019}
\instituicao{Universidade Federal de Santa Catarina}
\programa{Programa de Pós-Graduação em Ciências da Computação}
\local{Florianópolis} % Apenas cidade! Sem estado
\tese % ou \dissertacao
\titulode{Doutor em Ciência da Computação}
\orientador{Prof. Dr. Ben Trovato}
\coorientador{Prof. Dr. Lars Thørväld}
\centro{Centro Tecnológico}

\membrobanca{Prof. Valerie Béranger, Dr.}{Universidade Federal de Santa Catarina}
\membrobanca{Prof. Mordecai Malignatus, Dr.}{Universidade Federal de Santa Catarina}
\membrobanca{Prof. Huifen Chan, Dr.}{Universidade Federal de Santa Catarina}
\coordenador{Prof. Dr. Charles Palmer}
  \end{minted}
  \fonte{o autor.}
\end{figure}

Os elementos pré-textuais permanecem em grande parte tendo seu comportamento
determinado pelo \abnTeX. Alguns elementos incluídos nesse guia foram alterados
para satisfazer normas da UFSC. Esses elementos também presentes
nesse guia são obtidos com o código mostrado na \autoref{fig:pre}. Capa
(\mt|\imprimircapa|), folha de rosto (\mt|\imprimirfolhaderosto*|) listas
(\emph{e.g.} \mt|\listoffigures*|) e o sumário \mt|\tableofcontents*| são
obtidos através dos mesmos comandos do \abnTeX, mas agora respeitam novas
regras definidas pela BU.

A ficha catalográfica deve ser gerada em sistema próprio da BU, que produz um
PDF. Esse PDF deve ser salvo no seu projeto e incluído no documento.

A macro \mt|\incluirfichacatalografica{<pdf>}| faz essa inclusão usando o
pacote \href{https://www.ctan.org/pkg/pdfpages}{\mt|pdfpages|}. O \mt|\protect|
é necessário para evitar um \emph{bug} (uma linha em branco substitui o
\mt|\protect|).  A folha de certificação é gerada pelo comando
\mt|\imprimirfolhadecertificacao| de acordo (e na ordem) com os comandos de
dados emitidos anteriormente.

\begin{figure}[tb]
  \centering
  \caption{Elementos pré-textuais.}
  \label{fig:pre}
  \begin{minted}{latex}
\pretextual%
\imprimircapa%|*
\protect\incluirfichacatalografica{ficha.pdf}
\cleardoublepage \imprimirfolhadecertificacao
\cleardouplepage \listoffigures*
\cleardoublepage \tableofcontents*
  \end{minted}
  \fonte{o autor.}
\end{figure}

Antes de iniciar o primeiro capítulo é necessário chamar a macro \mt|\textual|,
como em todo documento \abnTeX. Após os elementos textuais, os pós-textuais
devem ser precedidos por um \mt|\postextual|. Para fazer uso da bibliografia,
basta fazer como mostra a \autoref{fig:pos}. A classe já inclui o
\texttt{abntex2cite} e o configura. A UFSC impôs poucas alterações no
tratamento de bibliografia.

\begin{figure}[tb]
  \centering
  \caption{Elementos pós-textuais.}
  \label{fig:pos}
  \begin{minted}{latex}
\postextual
\bibliography{example}
  \end{minted}
  \fonte{o autor.}
\end{figure}


\chapter{Problemas conhecidos}
\label{ch:problems}

\textbf{Atenção!} Há regras de formatação que não podem ser garantidas por essa
classe (ainda) devido a detalhes técnicos ou devido à forma como o \LaTeX e
BibTeX funcionam.

\section{Ordem de Elementos pré- e pós-textuais}
As regras da BU não alteram (nem alteravam antes) o significado e a ordem de
elementos pré-textuais definidos pela ABNT. Essa classe não tem como ajudá-lo.

\section{Alinhamento de Figuras}
O alinhamento de Figuras (ainda) não está totalmente claro nas regras da BU. Os
documentos da BU apresentam duas formas distintas. O autor do documento é
encorajado a continuar usando o alinhamento que usava até então. Essa dúvida
deverá ser esclarecida pela BU em breve, ocasião na qual essa classe será
atualizada.

\section{URLs de referências em inglês}

Alguns alunos costumam trocar os dizeres ``Disponível em'' gerados pelo
\texttt{abntex2cite} por ``Available at'' quando redigem o texto em inglês. No
entanto \textbf{isso é errado}! De acordo com o suporte da BU, apenas o corpo do
texto está em inglês, e assim como vários elementos pré-textuais permanecem em
português, os dizeres Disponível em deveriam permanecer em português.

Caso você deseje \textbf{contrariar a orientação acima}, a forma mais segura (no
sentido do LaTeX, não no sentido da entrega na BU) de \textbf{violá-la} é usar o
campo \texttt{note} do BibTeX. Você deverá fazer o equivalente a
\mt|Available at: \url{https://doi.org/10.1145/3338112}| em cada entrada no seu
arquivo \texttt{.bib}, além de desativar o campo \texttt{url} removendo-o ou renomeando-o
para. \texttt{x-url}.

\chapter{Referência de opções e macros}
\label{ch:ref}

\section{Opções da classe}

A classe possui apenas uma opção que fornece recursos: a primeira na lista
abaixo. Todas as demais servem para desativar comportamentos padrão.

\paragraph*{\texttt{english}} Prepara o \abnTeX{} e a classe para que o corpo
do documento esteja em inglês. Alguns elementos em português, especialmente os
relacionados a BU e a UFSC permanecem em português.

\paragraph*{\texttt{brazil}} Não tem efeito pois o idioma por padrão é o
português do Brasil. Entretanto, note que para utilizar outros idiomas além de
inglês e português, é necessário repassar todos os idiomas explicitamente à
classe \abnTeX.

\paragraph*{\texttt{embeddedlogo}} Se ativada, essa opção desempacota o logo da
UFSC diretamente de um base64 dentro do arquivo \texttt{.cls}. Isso permite que
o \texttt{.cls} seja um arquivo auto-contido. Pessoas avessas a \texttt{git
  submodule} e a diretórios podem simplesmente levar o \texttt{.cls} de um
projeto para outro. Entretanto, há uma desvantagem, e por isso essa opção não é
padrão. É necessário compilar com a opção \texttt{-shell-escape} em uma ambiente
\emph{UNIX-like}. O logo foi obtido da \href{http://identidade.ufsc.br/}{página
  de identidade visual da UFSC}.

\paragraph*{\texttt{times}, \texttt{lmodern}, \texttt{arial}} Escolhe a fonte
do documento. O padrão é \texttt{times}. No caso das escolhas \texttt{times} e
\texttt{arial} são usados clones \emph{open-source} (sem limitações legais à
redistribuição) que podem ser encontradas em qualquer instalação \LaTeX. A
opção \texttt{lmodern} corresponde à fonte Latin Modern, a fonte padrão do
\LaTeX. Essa é uma fonte serifada similar à Times, mas não possui respaldo
oficial da BU.

\paragraph*{\texttt{nocapautoref}} Essa classe configura o comando
\mt{\autoref} para que nomes de capítulos e seções sejam capitalizados. Ao
fornecer essa opção essa alteração deixará de ser feita.

\paragraph*{\texttt{noabntexcite}} Essa classe inclui o pacote
\texttt{abntex2cite} e o configura automaticamente. Caso seja necessário fazer
essa inclusão manualmente para fazer sua própria configuração, use essa opção
para que a classe não carregue o pacote.

\paragraph*{\texttt{nohidelinks}} O
\href{https://ctan.org/pkg/hyperref}{\texttt{hyperref}} por padrão desenha
bordas coloridas ao redor de \emph{hyperlinks}. Por padrão a classe desativa
esse comportamento, passando a opção \texttt{hidelinks} ao \texttt{hyperref}.
Ao ativar essa opção o \texttt{hyperref} voltará ao seu próprio comportamento
padrão.

\paragraph*{\texttt{noplainurl}} A norma ABNT NBR6023:2018 não utiliza mais URLs
na forma <http://example.org>, e apresenta a URL diretamente, sem os <>.  Esse
comportamento sempre é aplicado (a menos que seja usada a opção
\texttt{noabntexcite}). Para uniformizar o estilo, essa classe aplica a mesma
configuração no corpo do documento (nas instâncias onde se usaria
\mt|\url{http://example.org}|. Ao fornecer a opção \texttt{noplainurl} as
configurações do comando \mt|\url| não são alteradas.


\paragraph*{\texttt{nocleanheader}} O \abnTeX{} gera cabeçalhos na página que
incluem o título da seção. Essa classe suprime esse estilo de página,
proveniente da classe base \texttt{memoir}, pois eles não existem no template
da BU. Ao passar essa opção o cabeçalho padrão do \abnTeX{} será restaurado.

\paragraph*{\texttt{\textbackslash{}nocleardoublepage}} Por padrão, os comandos
\mt|\tableofcontents| \mt|\listoftables|, \mt|\listoffigures|,
\mt|\listoflistings| (se estiver definido) são modificados logo após o
\mt|\begin{document}| para que realizem um \mt|\cleardoublepage|. Logo,
esseselementos serão impressos garantidamente em uma nóva página que ficará no
anverso. Essa opção, se presente, desativa a injeção dessa alteração

\paragraph*{\texttt{*}} Quaisquer outras opções não listadas anteriormente são
repassadas à classe \abnTeX{} sem manipulação.

\section{Novas Macros}

Essa seção apresenta macros novas, não presentes no \abnTeX{} ou
\texttt{abntex2cite}.

\subsection{Comandos de dados}
\label{sec:dados}

Esses comandos visam fornecer dados que são posteriormente usados pelos
comandos de saída discutidos na próxima seção. \textbf{Importante:} não esqueça
de preencher os comandos de dados do próprio \abnTeX, como mostrado na
\autoref{sec:quick}.

\paragraph*{\texttt{\textbackslash{}programa\{<nome>\}}} Configura o nome do
Programa de Pós-Graduação.

\paragraph*{\texttt{\textbackslash{}titulode\{<texto>\}}} Define o título alvo.
O título \textit{deve ser escrito incluindo a especialidade}. Por exemplo,
\emph{Mestre em Ciência da Computação}.

\paragraph*{\texttt{\textbackslash{}tese}} Comando sem argumentos. Marca o
documento como uma tese de doutorado. Isso irá alterar o preâmbulo e a folha de
certificação.

\paragraph*{\texttt{\textbackslash{}dissertacao}} Comando sem argumentos. Marca
o documento como uma dissertação de mestrado. Esse é o padrão quando nem
\mt|\dissertacao| nem \mt|\tese| são usados.

\paragraph*{\texttt{\textbackslash{}nivel\{<nivel>\}}} Deve ser o nível do
título almejado, em minúsculas. O valor padrão é inferido a partir do uso (ou
não uso) de \mt|\tese| e \mt|\dissertacao|.

\paragraph*{\texttt{\textbackslash{}centro\{<texto>\}}} Centro da UFSC onde
está sediado o Programa de Pós-Graduação.

\paragraph*{\texttt{\textbackslash{}data\{<data por extenso>\}}} Define a data
a ser usada na folha de certificação. O ano dessa data será usado na capa e na
folha de rosto. Se esse comando for omitido, a data atual do sistema será
utilizada. Não utilize \mt|\today| como argumento em documentos em inglês. Isso
fará com que a data da folha de aprovação (que é em português) apareça em
inglês! Exemplo: \mt|\data{1 de Agosto de 2019}|.

\paragraph*{\texttt{\textbackslash{}ano\{<ano>\}} \\
            \texttt{\textbackslash{}mes\{<nome do mes>\}} \\
            \texttt{\textbackslash{}dia\{<dia>\}}}
Define cada componente da data individualmente.

\paragraph*{\texttt{\textbackslash{}membrobanca\{<nome com
títulos>\}\{<universidade por extenso>\}}}
Nome do professor avaliador da banca, incluíndo Prof. (ou Profa.) antes do
nome e o título \textbf{após} o nome. Exemplo: \emph{Prof. Mordecai
Malignatus, Dr.}. O orientador não deve ser incluído nessa seção. Não está
claro nas instruções da UFSC se o coorientador também não deveria. Note que
esse é o único lugar onde o título deve vir após o nome. Para orientador e
coordenador a BU coloca o título antes.

\paragraph*{\texttt{\textbackslash{}coordenadora\{<nome>\}}} Nome do
coordenador do programa. Nesse caso a BU usa títulos antes do nome. No caso de
coordenadora, o comando \mt|\coordenadora| pode ser usado. Exemplo:
\mt|\coordenador{ Prof. Dr. Mordecai Malignatus}|.

\subsection{Comandos de saída}

Esses comandos produzem resultados visíveis. Devem ser usados dentro do
ambiente \texttt{document}.

\paragraph*{\texttt{\textbackslash{}incluirfichacatalografica\{<arquivo>\}}}
Inclui a ficha catalográfica proveniente de \texttt{arquivo} que foi gerado em
\href{http://ficha.bu.ufsc.br/}{ficha.bu.ufsc.br}.

\paragraph*{\texttt{\textbackslash{}imprimirfolhadecertificacao}\\
            \texttt{\textbackslash{}imprimirfolhadeaprovacao}
} 
Imprime a ``folha de certificação'', que substitui a antiga ``folha de
aprovação''. Ambos os comandos produzem o mesmo resultado. Nessa folha
os membros da banca serão listados na ordem em que foram adicionados
com o comando \mt|\membrobanca{<nome>}{<universidade>}|.

\subsection{Comandos \abnTeX{} modificados}

\paragraph*{\texttt{\textbackslash{}imprimircapa}} O modo de uso
desse comando permanece o mesmo. No entanto, como as novas capas devem incluir o
brasão da UFSC, há alguns detalhes a serem observados. Ao utilizar a opção
\texttt{embeddedlogo}, o logo da UFSC é extraíde de um base64 embutido dentro do
próprio \texttt{.cls}. Sem usar essa opção é necessário que o arquivo
\texttt{logo-ufsc.pdf} seja encontrado pelo comando \mt|\includegraphics|. Para
que isso funcione, o \texttt{.cls} inicializa o \mt|\graphicspath| da seguinte forma:

\begin{minted}{latex}
  \graphicspath{%
    {.}%
    {ufsc-thesis-rn46-2019/}%
    {../ufsc-thesis-rn46-2019/}%
    {ufsc-thesis/}%
    {../ufsc-thesis/}%
\end{minted}

Essa inicialização é suficiente para os casos de uso típicos. Note que os
caminhos são relativos ao \emph{working directory} do \texttt{pdflatex}, não são
relativos ao \texttt{.cls} ou ao seu \texttt{.tex}.

\paragraph*{\texttt{\textbackslash{}imprimirfolhaderosto}} A folha de rosto foi
modifica de acordo com as intruções da BU. O elementa com maior diferênça é o
texto do preambulo (que é inserido abaixo do título). Nessa classe, o texto é
gerado automáticamente a partir da informação provida pelo uso de \mt|\tese| (ou
\mt|\dissertacao|), \mt|\programa| e \mt|\titulode|, apresentados
anteriormente. Ainda é possível usar \mt|\preambulo| diretamente para
sobreescrever o texto gerado. No entanto, leia as regras da BU com atenção.

\paragraph*{\texttt{\textbackslash{}begin\{dedicatoria\}...\textbackslash{}end\{dedicatoria\}}}
O ambiente \texttt{dedicatoria} do \abnTeX não aplica nenhuma formatação. O
ambiente foi alterado para que o texto fique alinhado a direita, com recuo de
5cm, ancorado na parte inferior da página. Essa formatação se originou do
documento da BU que explica as normas ABNT. A medida de 5cm foi tomada do
template \texttt{.doc} disponibilizado pela BU.

\paragraph*{\texttt{\textbackslash{}begin\{epigrafe\}...\textbackslash{}end\{epigrafe\}}}
O ambiente \texttt{epigrafe} do \abnTeX não aplica nenhuma formatação. O
ambiente foi alterado para que o texto fique alinhado a direita, com recuo de
4cm, espaçamento entre linhas simples e fonte 10. Essas são as mesmas
configurações do ambiente citação. o template \texttt{.doc} aplica todas, exceto
pela fonte 11, apesar de citar a ABNT NBR 10520. Visando maior consistência,
optou-se por usar a fonte 10pt.

\paragraph*{\texttt{\textbackslash{}begin\{abstract\}[<nome>]...\textbackslash{}end\{abstract\}}}
Em abnTeX2, não se deve chamar o ambiente abstract diretamente. Caso o alune o
chame, será interpretado como se o aluno tivesse chamado
\mt|\begin{resumo}[Abstract]| e a linguagem (controlando hifenização) será
alterada para inglês.

\paragraph*{\texttt{\textbackslash{}begin\{resumo\}[<nome>]...\textbackslash{}end\{resumo\}}}
Ambiente usado para gerar os resumos (para documentos em inglês é necessário um
resumo estendido) e \emph{abstracts}. A BU fornece além do template geral .doc,
fornece um template para o resumo estendido, que também deve ser gerado usando
esse ambiente. Nos três tipos de resumo exigidos, o aluno é responsável por
incluir um parágrafo ao final contendo as Palavras-chave\footnote{Sugestão:
  \mft|\vspace{\baselineskip}Palavras-chave: Kw1. Kw2. Kw3.|}: Regras
implementadas nesse ambiente:
\begin{itemize}
\item Espaçamento simples, fonte 12, sem indentação de parágrafos;
\item A formatação de \mt|\section| torna-se igual à formatação de
  \mt|\subsection|, para garantir cumprimento da formatação usada no template de
  resumo estendido;
\item Espaçamentos após títulos não são aplicados (resumo estendido);
\item A hifenização é alterada para \textit{brazil}, a menos que o autor tenha
  usado o ambiente abstract ou tenha usado Abstract, abstract ou ABSTRACT como
  argumento \texttt{<nome>} do ambiente.
\end{itemize}

\chapter{Mais exemplos de formatação}
\label{ch:ex}

Essa frase é verdadeira pois tem um \mt|\cite| no final \cite{turing1937}. Essa
é mais verdadeira ainda pois tem um (\mt|\cite{turing1937,dijkstra1968}|) no
final \cite{turing1937,dijkstra1968}. Já esta frase inofensiva usa
\mt|\citeonline{dijkstra1968}| para citar \citeonline{dijkstra1968}
nominalmente. O trabalho de \citeonline{diffie1976} foi altamente influente
\cite{diffie1976}. Essa outra frase cita o trabalho que \citeonline{Saleem2018}
escreveu com outros 4 autores. Para algo completamente novo, veja um footnote
com url\footnote{\url{http://example.org/}}

Mais algumas citações de tipos específicos de documentos:
\begin{itemize}
\item \mt|@inproceedings|: \citeonline{Ullman1989magic}. Jabuti
  \cite{Ullman1989magic}.
\item \mt|@article|: \citeonline{Distefano2019}, framboesa \cite{Distefano2019};
\item \mt|@book|: \citeonline{Abiteboul1995}, goiaba \cite{Abiteboul1995};
\item \mt|@incollection|: \citeonline{Forgy1989}, melancia \cite{Forgy1989};
\item \mt|@techreport|: \citeonline{rdf11}, figo \cite{rdf11}.
\end{itemize}

A lista abaixo mostra o efeito de \mt|\autoref{}| com capítulos e (sub)seções.

\begin{itemize}
\item Há coisas na \autoref{ch:ex};
\item Há coisas na \autoref{sec:stuff};
\item Há coisas na \autoref{sec:more};
\item Há coisas na \autoref{sec:yet-more};
\item Há coisas na \autoref{sec:yet-another} (\abnTeX{} come um ``sub''
  intencionalmente no português e o mantem em inglês).
\end{itemize}

Citações são feitas com \mt|\begin{citacao}...\end{citacao}|. A BU faz
as mesmas exigências que já são o \textit{default} na classe
\abnTeX\footnote{O alinhamento e o filete de notas de rodapé também
não necessitou de modificações, além do tamanho da fonte. Essa frase
não serve a nenhum propósito além de causar uma quebra de linha para
que o alinhamento seja avaliado.}.

\begin{citacao}
  A elaboração do trabalho de conclusão de curso em nível de mestrado
  e de doutorado na UFSC deverá atender aos critérios e procedimentos
  estabelecidos nesta resolução normativa e em diretrizes
  estabelecidas pela Pró-Reitoria de Pós-Graduação e pelos Programas
  de Pós-Graduação.
\end{citacao}

Atenção! O template da BU deixa figuras e tabelas alinhadas à esquerda. No
entanto, o tutorial de Word disponibilizado pela BU diz que o \mt|\caption{}|
(legenda acima) e o \mt|\fonte{}| (o ``Fonte: o autor.'', abaixo) devem
respeitar o ``alinhamento da ilustração''. Esse tutorial apresenta apenas uma
ilustração alinhada à esquerda). O tutorial explicando a ABNT mostra uma figura
centralizada com legendas alinhadas a esquerda e com recuo até o começo da
figura. O autor do \texttt{.cls} se exime de qualquer culpa. Alinhe aqui (com
\mt|\centering|, \mt|\flushright| ou \mt|\flushleft|) como mandar o seu
coração. Veja na \autoref{fig:logo} o efeito de se usar \mt|\centering|.

\begin{figure}[t]
  \centering
  \caption{Logotipo da Universidade Federal de Santa Catarina.}
  \label{fig:logo}

  \includegraphics[width=.2\linewidth]{../logo-ufsc.pdf}
  \fonte{o autor.}
\end{figure}

\section{Coisas}
\label{sec:stuff}
Imagine alguma afirmação de alto valor científico aqui.

\subsection{Outras coisas}
\label{sec:more}
Olá! Eu vim do passado para te avisar que o texto de uma dissertação deve ser
impessoal e você não deveria tentar conversar com o leitor.

\subsubsection{Outras coisas mais}
\label{sec:yet-more}
Estudos demonstram que essa afirmação é falsa.

\subsubsubsection{Ainda outras coisas mais}
\label{sec:yet-another}
Fazer a grama verde, como? Novamente o jogo foi perdido. Opcionalmente, tudo
pode ser opcional. Recursos foram gastos com isso. Descubra a verdade nas
capitalizadas.
% Fiquei 15 minutos mais próximo da morte ao escrever isso. Você pode chegar
% ainda mais perto se tentar entender.

%%%%%%%%%%%%%%%%%%%%%%%%%%%%%%%%%%%%%%%%%%%%%%%%%%%%%%%%%%%%%%%%%%%%
%%% Elementos pós-textuais                                       %%%
%%%%%%%%%%%%%%%%%%%%%%%%%%%%%%%%%%%%%%%%%%%%%%%%%%%%%%%%%%%%%%%%%%%%

\postextual
\bibliography{userguide}

\end{document}

% LocalWords:  pdf ufsc-thesis-rn abntex PPGCC Zambonin pull request GitHub PRs
% LocalWords:  alexishuf Issues alpha drop-in template cls shell-escape ZIP Git
% LocalWords:  UNIX-like Makefile git submodule shame on you Overleaf english
% LocalWords:  overview Pós-Graduação Prof BU pdfpages bug BibTeX URLs at bst
% LocalWords:  Available cite-alf url x-url brazil embeddedlogo auto-contido pt
% LocalWords:  lmodern arial open-source Latin Modern nocapautoref ufsc-thesis
% LocalWords:  working directory pdflatex tex dijkstra diffie Saleem Ullman rdf
% LocalWords:  magic Distefano Abiteboul Forgy default Word captions userguide
% LocalWords:  noplainurl noabntexcite begin dedicatoria end doc abstracts
% LocalWords:  corretude indentação
